% Generated draft: requires mathematical and source review.
\documentclass[11pt]{article}
\usepackage[margin=25mm]{geometry}
\usepackage{fontspec}
\setmainfont{Latin Modern Roman}
\usepackage{amsmath,amssymb,mathtools,unicode-math}
\setmathfont{Latin Modern Math}
\usepackage{enumitem,needspace,xurl,hyperref}
\setlength{\parindent}{0pt}
\setlength{\parskip}{5pt}
\setlength{\emergencystretch}{3em}
\newcommand{\N}{\mathbb{N}}
\newcommand{\Z}{\mathbb{Z}}
\newcommand{\Q}{\mathbb{Q}}
\newcommand{\R}{\mathbb{R}}
\newcommand{\C}{\mathbb{C}}
\newcommand{\F}{\mathbb{F}}
\newcommand{\A}{\mathbb{A}}
\newcommand{\PP}{\mathbb{P}}
\newcommand{\E}{\mathbb{E}}
\newcommand{\eps}{\varepsilon}
\newcommand{\dd}{\,\mathrm{d}}
\begin{document}
\section*{UnsolvedMath ID: 6200099}
\textbf{Source problem code:} catalogue entry
\subsection{Definitions and mathematical statement}
Let \(\mathcal T(S)\) be the Teichm\"uller space of a surface \(S\), equipped with the Teichm\"uller metric
\[
 d_T(X,Y)=\frac12\inf_F\log K(F),
\]
where \(F\) ranges over quasiconformal maps in the prescribed marking class and \(K(F)\) is maximal dilatation. The source asks whether this metric space has a quadratic isoperimetric inequality.

A standard continuous filling formulation is the following. For a Lipschitz loop \(\gamma:S^1\to\mathcal T(S)\), let \(L_T(\gamma)\) be its length and set
\[
 \operatorname{Fill}_T(\gamma)=\inf\{\operatorname{Area}_T(u):
 u:\overline{D^2}\to\mathcal T(S)\text{ is an admissible filling of }\gamma\}.
\]
With a specified parametrized area and filling regularity, ask whether there is \(C=C(S)<\infty\) such that
\[
 \operatorname{Fill}_T(\gamma)\le C L_T(\gamma)^2
\]
for every eligible loop. This states a quadratic upper bound; it does not assert a matching lower bound. The excerpt does not specify the surface type, Lipschitz versus other filling classes, metric-area convention, or continuous versus coarse/combinatorial interpretation. A coarse version may allow an additive constant or only large loops, and should not be identified silently with the displayed all-scale version.

The metric is explicitly the Teichm\"uller metric, not another metric on the same space. The summary reports quadratic Dehn-function results for mapping class groups and includes a speculative CAT(0)-comparison fragment, without precise citations. Neither is taken to establish the requested inequality for the full Teichm\"uller space. Its partial-progress label and the 2005 problem remain source assessments with no independently verified resolution here.
\subsection{Short English statement}
Can every loop in Teichm\"uller space, with its Teichm\"uller metric, be filled with area bounded by a constant times the square of its length? The surface and filling/area conventions need specification; related mapping-class-group results do not by themselves answer this metric-space question.
\subsection{Sources}
\begin{itemize}
\item Local registry: \url{data/unsolvedmath/problems.json}, permanent UnsolvedMath ID \(6200099\), source code catalogue entry; supplied statement, background, and research summary, updated 2026-08-06.
\item Misha Kapovich, \emph{Problems on Boundaries of Groups and Kleinian Groups} (2005), Problem 99, PDF page 24: \url{https://www.math.ucdavis.edu/~kapovich/EPR/problems.pdf}. Supplied normalized PDF-text extraction accessed 2026-07-29; current status marked as needing review.
\item The summary mentions mapping-class-group Dehn functions without identifiable authors, publication details, or a theorem for the stated Teichm\"uller-space inequality.
\end{itemize}
\end{document}
